\bookchapter{第二十九章：數字健康——和手機建立健康關係}

\chapterintro{你有沒有這樣的體驗：本來只想刷5分鐘短影片，結果一抬頭已經過去了兩小時？或者晚上躺在床上刷手機，越刷越精神，最後凌晨兩點才睡著？手機改變了我們的生活，但也在悄悄吞噬我們的時間和健康。這一章不是讓你“戒掉手機”，而是幫你重新奪回對手機的主動權。}

\booksection{一、先看看你的“螢幕時間”——數字讓人清醒}

在改變之前，先搞清楚現狀。手機裡有一個很多人從未開啟過的功能，叫“螢幕使用時間”（iPhone）或“數字健康”（安卓），它能告訴你每天到底在手機上花了多少時間，以及這些時間花在了哪些App上。

\begin{itemize}
  \item \textbf{iPhone}：設定 > 螢幕使用時間 > 檢視所有活動。你會看到每天的解鎖次數、拿起手機的次數、每個App的使用時長。
  \item \textbf{安卓}：設定 > 數字健康與家長控制。華為叫"健康使用手機"，小米叫"螢幕時間管理"。功能類似，都能看到每日使用資料。
\end{itemize}

\textbf{第一次看這個資料的人，80\%會被嚇到}。我有個朋友看完之後發現自己每天在抖音上花了4個多小時，一週就是28小時——相當於一份兼職的時間。這種“數字衝擊”是改變的第一步。

\booksection{二、給App設定“使用限額”——給自己裝個限速器}

知道問題在哪之後，就可以有針對性地設定限制。不需要一刀切把App刪了，只需要給它設個“每日額度”：

\begin{enumerate}
  \item \textbf{iPhone}：設定 > 螢幕使用時間 > App限額 > 新增限額。選中你要限制的App（比如抖音），設定每天30分鐘。到時間後App圖示會變灰，開啟需要額外輸入密碼。
  \item \textbf{安卓}：設定 > 數字健康 > 應用計時器。選中App後設定每日使用上限，到時間後App圖示會變灰。
  \item \textbf{小技巧}：把限額密碼設成一個你不常用的、不太好記的密碼。太容易輸的話，你會忍不住每次都"續命"，限額就形同虛設了。
\end{enumerate}

\booksection{三、“勿擾模式”的正確用法——不是關掉通知，是關掉焦慮}

很多人的手機一天響幾十次通知，每響一次注意力就被打斷一次。你以為你只是“看一眼通知”，實際上每次打斷後重新集中注意力需要15-20分鐘。勿擾模式不是讓你“失聯”，而是幫你重新掌控注意力。

\begin{itemize}
  \item \textbf{定時開啟}：設定勿擾模式在晚上10點到早上7點自動開啟。這樣睡覺時手機不會因為一條垃圾簡訊把你吵醒。
  \item \textbf{允許重要聯絡穿透}：在勿擾模式設定裡，把"允許來電"設為"收藏聯絡人"或"重複來電"（3分鐘內第二次來電會正常響鈴）。這樣你不會錯過家人的緊急電話，但也不會被推銷電話打擾。
  \item \textbf{工作專注模式}：iPhone的"專注模式"可以設自定義場景——比如"工作"模式下只允許工作相關App通知，"閱讀"模式下所有通知全靜音。
\end{itemize}

\booksection{四、護眼不靠App靠習慣——20-20-20法則}

長時間看手機螢幕，眼睛乾澀、酸脹是必然的。護眼模式（藍光過濾）有一定幫助，但真正有效的是改變用眼習慣：

\begin{tcbinfo}{20-20-20法則}
每看螢幕\textbf{20分鐘}，把視線移到\textbf{20英尺（約6米）}以外的地方，保持\textbf{20秒}。這個簡單的動作能有效緩解睫狀肌疲勞。

\end{tcbinfo}

\begin{itemize}
  \item \textbf{開啟護眼模式}：設定 > 顯示 > 護眼模式（或夜覽模式），開啟後螢幕色溫變暖，減少藍光。建議設為日落到日出自動開啟。
  \item \textbf{調整螢幕亮度}：環境暗時螢幕也要暗。很多人在黑暗中看手機亮度調太高，這是傷眼最嚴重的行為。開啟"自動亮度"並確保它正常工作。
  \item \textbf{保持距離}：手機距離眼睛至少30釐米。很多人喜歡把手機拿得特別近看，這樣睫狀肌一直在緊張收縮狀態，近視加深很快。
\end{itemize}

\booksection{五、“低頭族”的頸椎自救指南}

低頭看手機時，你的脖子承受的壓力是正常站立時的5-6倍。長期低頭會導致頸椎曲度變直、骨質增生、頭暈手麻。以下是幾個簡單的改善方法：

\begin{itemize}
  \item \textbf{把手機舉高}：看手機時儘量把手機舉到與眼睛平齊的高度，而不是低頭看放在腿上的手機。雖然手臂會累，但脖子會感謝你。
  \item \textbf{每30分鐘活動一下脖子}：慢慢向左轉頭到極限，停3秒，回正；再向右轉，停3秒。重複5次。然後做"米字操"——用下巴在空中寫一個"米"字。
  \item \textbf{靠牆站立3分鐘}：後腦勺、肩胛骨、臀部、腳後跟四點貼牆，每天站3分鐘，能有效糾正含胸低頭的體態。
\end{itemize}

\booksection{六、睡眠與手機——臥室裡最該改的習慣}

手機螢幕的藍光會抑制褪黑素分泌，讓你“越看越精神”。很多人覺得“我刷著刷著就困了”，實際上是被藍光欺騙了大腦，等你放下手機反而不容易入睡。以下是改善建議：

\begin{itemize}
  \item \textbf{睡前1小時不看手機}：這是最有效也最難做到的一條。如果做不到，至少開"夜覽/護眼模式"把藍光降到最低。
  \item \textbf{手機放在伸手夠不到的地方}：把充電器放在臥室另一頭，而不是床頭櫃。這樣你既不會忍不住拿起來刷，鬧鐘響了你也得起床才能關掉——一舉兩得。
  \item \textbf{開啟"睡眠模式"}：iPhone的"睡眠"專注模式可以自動在睡眠時段鎖定螢幕、簡化鎖屏；安卓各品牌也有類似功能。
  \item \textbf{用鬧鐘代替手機}：買一個十幾塊錢的實體鬧鐘，不用手機當鬧鐘，你會發現少了"看一眼時間然後忍不住刷一下"的陷阱。
\end{itemize}

\booksection{七、資訊焦慮症——學會“選擇性失明”}

我們每天接收的資訊量，是古人一輩子接收的量。但人的大腦不是為處理這麼多資訊設計的，資訊過載會導致焦慮、注意力渙散、決策疲勞。解決方法不是“獲取更多資訊”，而是“主動遮蔽”：

\begin{itemize}
  \item \textbf{取消關注}：開啟你關注的所有公眾號、UP主、博主列表，問自己一個問題："過去一個月我從中獲得過實際價值嗎？"如果沒有，取消關注。資訊渠道貴精不貴多。
  \item \textbf{關閉非必要App推送}：除了微信、電話、簡訊、日曆之外，其他App的通知全關掉。新聞不需要推送——你主動去看的時候自然會看到。
  \item \textbf{設立"資訊齋戒日"}：每週選一天（比如週日），全天不看手機新聞、不刷短影片。你會發現天不會塌下來，而且週一精神狀態明顯更好。
\end{itemize}

\begin{tcbpractice}{動手：一週數字健康挑戰}
\begin{enumerate}
  \item 今天：開啟螢幕使用時間/數字健康，檢視你上週的使用資料。
  \item 明天：給耗時最長的娛樂App設定每日限額（建議30-60分鐘）。
  \item 後天：設定勿擾模式定時開啟（22:00-07:00），允許收藏聯絡人來電。
  \item 第四天：睡前1小時不看手機，手機充電器放到臥室另一頭。
  \item 第五天：取消關注5個你不再看的公眾號或博主。
  \item 第六天：嘗試一整天的"資訊齋戒"——不看新聞不刷短影片。
  \item 第七天：檢視本週的螢幕使用時間報告，和上週對比。
\end{enumerate}

\end{tcbpractice}

\begin{tcbsummary}{本章小結}
\begin{itemize}
  \item 先看螢幕使用時間資料，認清現狀是改變的第一步
  \item 給娛樂App設限額，給勿擾模式設定時，重新掌控注意力
  \item 20-20-20法則是最簡單有效的護眼習慣
  \item 低頭看手機時脖子承重5-6倍，舉高手機+定期活動脖子
  \item 睡前1小時不看手機，充電器放到伸手夠不到的地方
  \item 資訊焦慮的解藥是"主動遮蔽"，不是"獲取更多"
\end{itemize}

\end{tcbsummary}
